\subsection{GNSS Factor}
The GNSS factor represents a direct measurement constraint that links the estimated vehicle pose to the absolute position and heading obtained from the dual-antenna GNSS receiver. For AUVs, GNSS is mainly useful when the vehicle is at the surface, since the signal is not available when submerged.
\\ \\
The GNSS factor can be expressed as:
$$
    \mathbf{r}_{GNSS} = \mathbf{z}_{GNSS} - h(\mathbf{x})
$$
where $\mathbf{z}_{GNSS}$ is the measurement vector containing the absolute position and yaw angle, and $h(\mathbf{x})$ is the measurement model defined in Equation \ref{eq:aiding-measurement-model-gnss}. The measurement noise is modeled as:
$$
    \mathbf{r}_{GNSS} \sim \mathcal{N}(0, \mathbf{R}_{GNSS})
$$
where $\mathbf{R}_{GNSS}$ is the GNSS measurement covariance, as defined in Equation \ref{eq:aiding-measurement-model-gnss-noise}.
\\ \\
In the factor graph, the GNSS observation is implemented as a unary factor connected to the current pose node:
$$
    \| \mathbf{z}_{GNSS} - h(\mathbf{x}) \|^2_{\mathbf{R}_{GNSS}^{-1}}
$$
By incorporating both position and yaw, the GNSS factor provides a global reference when available. It helps anchor the estimated trajectory in the navigation frame and reduces long term drift during surface operation.